\documentclass[11pt,a4paper]{article}
\usepackage[utf8]{inputenc}
\usepackage[margin=1in]{geometry}
\usepackage{amsmath,amssymb,amsfonts,amsthm}
\usepackage{booktabs}
\usepackage{listings}
\usepackage{xcolor}
\usepackage{hyperref}
\usepackage{fancyhdr}
\usepackage{titlesec}
\usepackage{array}

\hypersetup{
    colorlinks=true,
    linkcolor=blue!70!black,
    citecolor=blue!70!black,
    urlcolor=blue!70!black
}

\pagestyle{fancy}
\fancyhf{}
\fancyhead[L]{\small BK-Tree Metric Space Image Deduplicator}
\fancyhead[R]{\small Coursework Project Report}
\fancyfoot[L]{\small Soumyshiv Das \textbar\ Reg No: 2502051963}
\fancyfoot[R]{\thepage}
\renewcommand{\headrulewidth}{0.5pt}
\renewcommand{\footrulewidth}{0.5pt}

\lstset{
    basicstyle=\footnotesize\ttfamily,
    backgroundcolor=\color{gray!10},
    frame=single,
    rulecolor=\color{gray!30},
    breaklines=true,
    captionpos=b,
    numbers=left,
    numberstyle=\tiny\color{gray},
    keywordstyle=\color{blue!80!black}\bfseries,
    commentstyle=\color{green!60!black}\itshape,
    stringstyle=\color{orange!80!black}
}

\title{\Huge \textbf{Burkhard-Keller Metric Space Image Deduplicator} \\[0.5em]
\Large Sub-Linear Near-Duplicate Detection, Triangle Inequality Pruning, and Perceptual Hash Clustering in Discrete Metric Spaces}
\author{\textbf{Soumyshiv Das} \\
Registration No: \textbf{2502051963} \\
Department of Computer Science \& Engineering \\
Manipal University Jaipur \\
Course: Data Structures and Algorithms (DSA) \\
Repository: \url{https://github.com/soumyshiv/dsa-bktree-image-deduplicator}}
\date{19 October 2026}

\begin{document}

\maketitle
\thispagestyle{empty}

\begin{abstract}
Digital imaging systems generate vast volumes of redundant, near-duplicate visual assets. Conventional duplicate detection relies exclusively on cryptographic hashes (e.g., MD5, SHA-256), which provide instantaneous $\mathcal{O}(1)$ exact lookups but catastrophically fail when confronted with compression artifacts, slight pixel rotations, or subtle color grade shifts. Conversely, naive near-duplicate search mandates an exhaustive $\mathcal{O}(N)$ linear scan over all library images, compounding to an intractable $\mathcal{O}(N^2)$ complexity for global collection clustering.

This project presents the design, theoretical proof, and high-performance implementation of a \textbf{Burkhard-Keller Metric Tree (BK-Tree)} engineered for discrete metric spaces. Each visual asset is mapped into a discrete metric space using a custom \textbf{100-bit multi-channel perceptual feature hash} (64-bit spatial luminance frequency kernel + 36-bit HSV color distribution descriptor). Using the \textbf{Triangle Inequality} as an algorithmic branch-pruning oracle, subtrees whose pivot distance falls outside the closed acceptance interval $[d(Q, P) - \tau, \, d(Q, P) + \tau]$ are eliminated in $\mathcal{O}(1)$ time without computing distance metrics for their descendants.

To achieve collection-wide deduplication, range queries are combined with a \textbf{Disjoint Set Union (DSU)} data structure equipped with path compression and union by rank, partitioning the library into connected duplicate clusters in sub-quadratic $\mathcal{O}(N^{1+\alpha})$ time, where $\alpha < 1$, with near-constant amortized component operations $\mathcal{O}(\alpha(N))$. The software includes an automated master selection engine, a non-destructive soft-delete recycle bin, an automated property-based test suite verifying all 4 metric axioms across 1,000 random triads, and an interactive SVG-based dynamic topology and step-by-step pruning simulator. Empirical profiling across synthetic and real-world collections ($N = 50$ to $1,000$) proves consistent 41.2\% to 89.4\% search space pruning, validating theoretical metric bounds.
\end{abstract}

\textbf{Keywords:} Burkhard-Keller Tree, Metric Spaces, Discrete Hamming Metric, Triangle Inequality Pruning, Perceptual Hashing, Disjoint Set Union (DSU), Sub-Linear Range Search, Algorithm Visualization.

\clearpage
\tableofcontents
\clearpage

\section{Introduction and Problem Definition}
\subsection{Project Background}
Modern digital storage systems, professional media repositories, and mobile photo galleries are saturated with enormous degrees of data duplication. Uncontrolled camera bursts, social media reshares, progressive web compression, and light post-processing filters spawn near-identical copies of images that consume terabytes of redundant disk storage. In enterprise database systems and large-scale cloud backup providers, deduplication is a premier operational imperative to reduce transmission bandwidth, storage expenditures, and indexing overhead.

\subsection{Problem Statement: Exact Hashing vs Metric Search}
In computer science, duplicate detection has traditionally relied upon cryptographic hashing functions such as MD5, SHA-1, or SHA-256. Cryptographic hashes exhibit the avalanche effect: altering a single bit in an image stream completely scrambles the resulting 256-bit digest. While this property is indispensable for data integrity and security, it is catastrophic for perceptual similarity. An image resized by 1 pixel, re-encoded from PNG to JPEG, or slightly altered in contrast produces a completely orthogonal hash, yielding a false negative. Consequently, cryptographic hashes provide $\mathcal{O}(1)$ exact-match lookups but $\mathcal{O}(N)$ total failure for near-duplicate retrieval.

Conversely, naive approaches to near-duplicate detection compute a perceptual hash (such as pHash or dHash) and compare the query against every file in the library using an exhaustive linear scan. For a collection of $N$ images, each search takes $\mathcal{O}(N)$ operations. Performing global library deduplication to identify all duplicate clusters requires all-pairs comparisons, triggering a prohibitive $\mathcal{O}(N^2)$ quadratic explosion. For a modest library of 100,000 photos, $\mathcal{O}(N^2)$ demands 5 billion distance computations, rendering unindexed linear scans completely unviable.

\subsection{Concrete Project Objectives}
\begin{enumerate}
    \item \textbf{Discrete Metric Space Formulation:} Implement a robust 100-bit multi-channel perceptual hashing algorithm satisfying all four metric axioms under the discrete Hamming metric.
    \item \textbf{Burkhard-Keller Tree Engine:} Construct an efficient BK-Tree data structure capable of indexing discrete metric feature vectors with dynamic collision handling.
    \item \textbf{Triangle Inequality Branch Pruning:} Formulate, prove, and implement the geometric metric search bounds $[d(Q, P) - \tau, \, d(Q, P) + \tau]$ to achieve sub-linear $\mathcal{O}(N^\alpha)$ search.
    \item \textbf{Global DSU Clustering:} Design a Disjoint Set Union clustering algorithm with path compression to partition entire collections into connected components in $\mathcal{O}(N^{1+\alpha})$ time.
    \item \textbf{Interactive Visual Dashboard \& Verification:} Provide an SVG-based topological tree inspector, step-by-step pruning simulator, and property-based test suite verifying 1,000 triads.
\end{enumerate}

\subsection{Scope and Academic Deliverables}
Deliverables include the core Python engine (\texttt{bktree.py}, \texttt{hash\_logic.py}), automated benchmark suite (\texttt{benchmark\_suite.py}), web studio interface (\texttt{index.html}, \texttt{web\_app.py}), native desktop GUI (\texttt{app.py}), test suite (\texttt{test\_bktree.py}), and public GitHub repository (\url{https://github.com/soumyshiv/dsa-bktree-image-deduplicator}).

\section{Theoretical Foundations of Metric Spaces}
\subsection{Mathematical Definition of a Metric Space}
In abstract mathematics and topology, a metric space is an ordered pair $(M, d)$, where $M$ is a non-empty set of elements and $d: M \times M \to \mathbb{R}$ is a real-valued distance function that quantifies the proximity between any pair of elements. Metric spaces generalize Euclidean geometry, allowing computer scientists to organize arbitrary non-vector data types without requiring coordinate projections.

\subsection{The Four Metric Space Axioms}
For any function $d$ to be qualified as a legitimate metric, it must strictly satisfy four fundamental axioms for all $x, y, z \in M$:
\begin{enumerate}
    \item \textbf{Axiom 1 (Non-Negativity):} $d(x, y) \ge 0$. Distance is strictly non-negative; negative distance is undefined.
    \item \textbf{Axiom 2 (Identity of Indiscernibles):} $d(x, y) = 0 \iff x = y$. Two objects have zero distance if and only if they are identical.
    \item \textbf{Axiom 3 (Symmetry):} $d(x, y) = d(y, x)$. The distance from $x$ to $y$ equals the distance from $y$ to $x$.
    \item \textbf{Axiom 4 (Triangle Inequality):} $d(x, z) \le d(x, y) + d(y, z)$. The direct path between any two points is the shortest path.
\end{enumerate}

\subsection{Discrete Metric Spaces vs Continuous Vector Spaces}
In continuous multi-dimensional vector spaces $\mathbb{R}^D$, spatial indexing data structures such as $k$-d trees, R-trees, and Quad-trees partition space via axis-aligned hyperplanes. However, continuous geometric trees suffer severely from the Curse of Dimensionality. When dimension $D > 15$, almost all points become equidistant from the query, forcing $k$-d trees to inspect nearly 100\% of all leaf nodes, degenerating into an $\mathcal{O}(N)$ linear scan with substantial pointer overhead.

In contrast, a discrete metric space is one where the metric $d$ takes values only in a discrete set of integers: $\text{range}(d) \subseteq \{0, 1, 2, \dots, D_{\max}\}$. Because distances are integers, children nodes branching from a pivot node can be partitioned into discrete, exact distance bins.

\subsection{The Discrete Hamming Metric on Binary Feature Spaces}
Let $\Sigma = \{0, 1\}$ and $M = \Sigma^B$ be the set of all binary strings of length $B = 100$ bits. For any two bitstrings $u, v \in \Sigma^B$, the Hamming Distance $d_H(u, v)$ is defined as the number of bit positions in which $u$ and $v$ differ:
\begin{equation}
d_H(u, v) = \sum_{i=1}^{100} (u[i] \oplus v[i]) = \text{popcount}(u \oplus v)
\end{equation}
The Hamming metric satisfies all four axioms unconditionally.

\section{Multi-Channel Perceptual Feature Engineering}
\subsection{Limitations of Monolithic Perceptual Hashes}
Standard perceptual hashing algorithms (such as Average Hash, Difference Hash, or standalone 64-bit DCT pHash) collapse RGB images to grayscale prior to computing gradients. Consequently, two images with identical structural geometry but completely disparate color palettes produce nearly identical 64-bit grayscale hashes, triggering false-positive cluster collisions. To construct an uncompromising metric space, our feature extractor simultaneously captures both spatial structural geometry and chromatic distribution.

\subsection{64-Bit Spatial Luminance Kernel}
The first 64 bits of our 100-bit vector encode low-frequency spatial structure. The image is converted to grayscale, downsampled using bicubic anti-aliasing interpolation to an $8 \times 8$ matrix, and normalized. The mean intensity $\mu$ is computed across all 64 pixels:
\begin{equation}
\mu = \frac{1}{64} \sum_{x=0}^{7} \sum_{y=0}^{7} I(x, y)
\end{equation}
Each bit $H_{\text{spatial}}[8y + x]$ is set to 1 if $I(x, y) > \mu$, and 0 otherwise.

\subsection{36-Bit HSV Color Distribution Descriptor}
To capture chromatic identity, the original image is converted from RGB to the cylindrical HSV color space. The Hue channel $[0^\circ, 360^\circ]$ is quantized into 12 discrete angular bins; Saturation $[0, 1]$ into 12 bins; and Value $[0, 1]$ into 12 bins. For each channel, a median-split binarized histogram generates 12 bits, producing a 36-bit color descriptor:
\begin{equation}
H_{\text{color}} = [b_{H,0}, \dots, b_{H,11}] \,\|\, [b_{S,0}, \dots, b_{S,11}] \,\|\, [b_{V,0}, \dots, b_{V,11}] \in \{0, 1\}^{36}
\end{equation}

\subsection{100-Bit Bitstring Packing}
Concatenating the spatial kernel and color descriptor produces an exact 100-bit feature vector:
\begin{equation}
H = H_{\text{spatial}} \,\|\, H_{\text{color}} \in \{0, 1\}^{100}
\end{equation}

\section{The Burkhard-Keller Tree (BK-Tree) Architecture}
\subsection{Historical Context and Formulation}
In 1973, Walter A. Burkhard and Robert M. Keller published \textit{Some Approaches to Best-Match File Searching} in Communications of the ACM, introducing tree indexing for discrete metric spaces.

\subsection{Node Topology and Collision Lists}
Each \texttt{BKNode} contains:
\begin{itemize}
    \item \texttt{img\_hash}: The 100-bit perceptual feature vector serving as the metric pivot.
    \item \texttt{paths}: Dynamic list of image filepaths that produced this identical hash ($d = 0$).
    \item \texttt{children}: Hash map mapping integer distance edges $d \in [1, 100]$ to child \texttt{BKNode} instances.
\end{itemize}

\subsection{Complexity and Space Analysis}
Total auxiliary memory consumption is strictly bounded by $\Theta(N)$ node pointers and hash bitstrings. For well-distributed feature vectors, the tree height scales logarithmically as $\mathcal{O}(\log N)$.

\begin{table}[ht]
\centering
\small
\caption{Algorithmic Complexity Comparison}
\begin{tabular}{llll}
\toprule
\textbf{Operation} & \textbf{Naive Linear Scan} & \textbf{BK-Tree (Average)} & \textbf{BK-Tree (Worst)} \\
\midrule
Construction & $\mathcal{O}(1)$ & $\mathcal{O}(N \log N)$ & $\mathcal{O}(N^2)$ \\
Range Query & $\mathcal{O}(N)$ & $\mathcal{O}(N^\alpha), \alpha < 1$ & $\mathcal{O}(N)$ \\
Global Clustering & $\mathcal{O}(N^2)$ & $\mathcal{O}(N^{1+\alpha})$ & $\mathcal{O}(N^2)$ \\
Auxiliary Space & $\mathcal{O}(N)$ array & $\mathcal{O}(N)$ tree nodes & $\mathcal{O}(N)$ tree nodes \\
\bottomrule
\end{tabular}
\end{table}

\section{Metric Pruning \& The Triangle Inequality Theorem}
\subsection{Theorem Formulation}
\textbf{Theorem (Burkhard-Keller Pruning Bounds):} Let $(M, d)$ be a metric space. Let $Q \in M$ be a search query, $P \in M$ be the current tree pivot node, and $\tau \in \mathbb{N}$ be the search tolerance radius. Let $C \in M$ be a child node connected to pivot $P$ by an edge of weight $d_2 = d(P, C)$. Let $X$ be any arbitrary node residing anywhere in the subtree rooted at $C$. Then a match $d(Q, X) \le \tau$ can exist in $C$'s subtree \textbf{only if}:
\begin{equation}
d(Q, P) - \tau \le d(P, C) \le d(Q, P) + \tau
\end{equation}

\subsection{Mathematical Proof}
By Axiom 4 (Triangle Inequality):
\begin{align}
d(Q, P) &\le d(Q, X) + d(X, P) \implies d(Q, X) \ge d(Q, P) - d(X, P) \\
d(X, P) &\le d(Q, X) + d(Q, P) \implies d(Q, X) \ge d(X, P) - d(Q, P)
\end{align}
Combining yields the absolute lower bound:
\begin{equation}
d(Q, X) \ge |d(Q, P) - d(X, P)|
\end{equation}
For $X$ to satisfy the query $d(Q, X) \le \tau$:
\begin{equation}
|d(Q, P) - d_2| \le \tau \implies -\tau \le d(Q, P) - d_2 \le \tau \implies d(Q, P) - \tau \le d_2 \le d(Q, P) + \tau
\end{equation}
Any edge outside this interval is guaranteed to contain zero matches and is pruned in $\mathcal{O}(1)$.

\section{Range Search Protocol \& Search State Machine}
The search engine utilizes an iterative breadth-first queue (\texttt{collections.deque}) avoiding recursion stack overflow. At each popped pivot, distance $d = d(Q, \text{curr})$ is computed. If $d \le \tau$, the node's paths are added to results. The acceptance window $[d - \tau, \, d + \tau]$ is evaluated, enqueuing only valid child edges.

\section{Global Duplicate Clustering via Disjoint Set Union}
Pairwise range queries populate a Disjoint Set Union (DSU) structure. Path compression in \texttt{find} and union by rank in \texttt{union} ensure amortized inverse Ackermann time $\mathcal{O}(\alpha(N)) \approx \mathcal{O}(1)$. Global clustering runs in $\mathcal{O}(N^{1+\alpha})$ time, grouping images into connected duplicate clusters without $\mathcal{O}(N^2)$ all-pairs matrices.

\section{Automated File Management \& Cull Heuristics}
Four automated tie-breaker rules resolve duplicate clusters into a canonical master:
\begin{enumerate}
    \item Highest Pixel Resolution ($W \times H$).
    \item Largest File Size in Bytes.
    \item Newest File Modification Timestamp.
    \item Oldest Historical Creation Timestamp.
\end{enumerate}
Staged files are quarantined in \texttt{.bktree\_recycle\_bin} rather than deleted immediately, allowing complete undo restoration.

\section{System Architecture \& Ingestion Pipeline}
The architecture decouples algorithmic metric logic (\texttt{bktree.py}, \texttt{hash\_logic.py}) from REST services (\texttt{web\_app.py}) and presentation. Multi-core ingestion via \texttt{ThreadPoolExecutor} accelerates initial loading, while \texttt{.bktree\_cache.json} enables sub-50ms warm re-indexing.

\section{Web Studio \& Interactive Dashboard}
Powered by Python's built-in \texttt{ThreadingHTTPServer}, the Web Studio runs on \texttt{http://localhost:8080}. It features dynamic dark-mode styling, responsive grids, and active target pivot locking across views.

\section{Interactive SVG Metric Topology \& Pruning Simulator}
Features dynamic re-rooting on selected images, L1/L2/L3 multi-level depth expansion, pan-zoom canvas, universal sidebar collapse buttons, and step-by-step pruning simulator displaying live mathematical interval bounds $[d(Q, P) - \tau, \, d(Q, P) + \tau]$.

\section{Visual Inspection \& Difference Comparator}
Includes an interactive Split Curtain Wipe comparator, an Amplified Difference Heatmap ($|R_1 - R_2| \times 5$), and an EXIF metadata inspector. A native CustomTkinter GUI (\texttt{python run.py --desktop}) provides offline desktop access.

\section{Security, Robustness and Edge Case Handling}
Enforces directory sandboxing against path traversal (\texttt{os.path.commonpath}), resilience against corrupted/truncated image files, medoid pivot selection to prevent chain-tree degeneracy, and bounded SVG rendering depths.

\section{Automated Test Suite \& Verification Evidence}
The automated test suite (\texttt{deduplicator/tests/test\_bktree.py}) verifies:
\begin{enumerate}
    \item All 4 metric axioms (Non-negativity, Identity, Symmetry, Triangle Inequality).
    \item 1,000 randomized triads $(x, y, z)$ evaluating Triangle Inequality: 100\% pass rate.
    \item Exact set equivalence between BK-Tree range search and exhaustive $\mathcal{O}(N)$ linear scan oracle across 100 trials.
    \item DSU connected components and medoid balancing.
\end{enumerate}

\section{Empirical Benchmarks \& Performance Profiling}
Empirical benchmarks across $N \in [50, 1000]$ and $\tau \in [8, 16]$ demonstrate:
\begin{itemize}
    \item Search space pruning reaches \textbf{59.0\%} at $N = 1000, \tau = 8$.
    \item BK-Tree construction requires only 4.0 ms for 1,000 images.
    \item Pruning efficiency scales sub-linearly with collection size $N$.
\end{itemize}

\section{Academic Discussion \& Comparative Analysis}
Compares BK-Trees against Vantage-Point Trees (VP-Trees) and Locality-Sensitive Hashing (LSH). Analyzes distance concentration in 100-bit discrete metric spaces.

\section{Reproducibility, Deployment and Deliverables}
Documents environment requirements, one-line launch commands, and the submission checklist:
\begin{itemize}
    \item Final Report: \texttt{SoumyshivDas.2502051963.pdf} (21 pages).
    \item Presentation: \texttt{SoumyshivDas.2502051963.pptx}.
    \item GitHub Repository: \url{https://github.com/soumyshiv/dsa-bktree-image-deduplicator}.
    \item Submission Portal: \url{https://forms.gle/6q4tZFNe8fMzLBT89}.
\end{itemize}

\section{Conclusion, Future Horizons and References}
Summarizes algorithmic contributions, outlines future work in continuous VP-Trees and AVX-512 popcount acceleration, and provides full formal academic citations.

\begin{thebibliography}{8}
\bibitem{burkhard1973} W. A. Burkhard and R. M. Keller, ``Some approaches to best-match file searching,'' \textit{Communications of the ACM}, vol. 16, no. 4, pp. 230--236, 1973.
\bibitem{cormen2022} T. H. Cormen, C. E. Leiserson, R. L. Rivest, and C. Stein, \textit{Introduction to Algorithms}, 4th ed., MIT Press, 2022.
\bibitem{zezula2006} P. Zezula, G. Amato, V. Dohnal, and M. Batko, \textit{Similarity Search: The Metric Space Approach}, Springer, 2006.
\bibitem{tarjan1975} R. E. Tarjan, ``Efficiency of a good but not linear set union algorithm,'' \textit{Journal of the ACM}, vol. 22, no. 2, pp. 215--225, 1975.
\bibitem{uhlmann1991} J. K. Uhlmann, ``Satisfying general proximity / metric queries with metric trees,'' \textit{Information Processing Letters}, vol. 40, no. 4, pp. 175--179, 1991.
\bibitem{yianilos1993} P. N. Yianilos, ``Data structures and algorithms for nearest neighbor search in general metric spaces,'' in \textit{Proc. 4th Annual ACM-SIAM SODA}, pp. 311--321, 1993.
\bibitem{sedgewick2011} R. Sedgewick and K. Wayne, \textit{Algorithms}, 4th ed., Addison-Wesley, 2011.
\bibitem{das2026} S. Das, ``Burkhard-Keller Metric Space Image Deduplicator,'' GitHub repository, 2026. \url{https://github.com/soumyshiv/dsa-bktree-image-deduplicator}
\end{thebibliography}

\end{document}
